\section{Report Contract}

Every W-series method emits the same report shape so that any analytic can be
reproduced from public inputs and read with its caveats attached. The per-record
transcript fields are:

\begin{center}
\small
\begin{tabularx}{\textwidth}{lX}
\toprule
Field & Meaning \\
\midrule
\texttt{method\_id} & Registered forensic analytic. \\
\texttt{feature\_set} & Variables used by the model. \\
\texttt{baseline\_population} & Comparison units and time window. \\
\texttt{model\_id} & Residual, digit test, spatial model, change-point, etc. \\
\texttt{unit\_id} & Reporting unit, batch, scanner, or contest scored. \\
\texttt{score} & Anomaly score or test statistic. \\
\texttt{p\_value\_or\_rank} & Optional statistical calibration. \\
\texttt{investigation\_note} & Machine-readable caveat for the scored unit. \\
\bottomrule
\end{tabularx}
\end{center}

A report is also bound to its inputs. Every W-series report must carry the
source package hashes it was computed from, the feature set and baseline
population that define the comparison, the model and method identity, the
unit-level scores, the caveats and false-positive boundary, and a plain
statement that the output is investigative and not certifying. Binding the
report to source package hashes is what makes it reproducible: a second analyst
with the same RCOUNT package and the same registered method should recompute the
same scores.

The report contract is specified here; its implementation placement is an open
decision. The analytics could live as RCOUNT audit-style reports beside the
existing transcripts, or in a later dedicated forensics crate. Either way, the
W-series transcript must remain a distinct artifact from V-series certification
evidence, so that an anomaly score is never stored or surfaced as a pass or
fail.
